\documentclass[12pt]{article}

\usepackage[paperwidth=200mm,paperheight=112.5mm,margin=11mm,top=10mm]{geometry}
\usepackage{amsmath,amssymb,booktabs,array,enumitem}
\usepackage[T1]{fontenc}
\usepackage{lmodern}

\pagestyle{empty}
\setlength{\parindent}{0pt}
\setlength{\abovedisplayskip}{9pt}
\setlength{\belowdisplayskip}{9pt}

\newcommand{\R}{\mathbb{R}}
\newcommand{\T}{\mathsf{T}}
\DeclareMathOperator{\col}{col}
\DeclareMathOperator{\sign}{sign}
\newcommand{\fbase}{f_{\mathrm{base}}}
\newcommand{\ftruth}{f_{\mathrm{truth}}}
\newcommand{\ueff}{u_{\mathrm{eff}}}
\newcommand{\cc}{c_{\mathrm{c}}}
\newcommand{\panel}[1]{\subsection*{\Large #1}\vspace{1mm}}
\newcolumntype{L}[1]{>{\raggedright\arraybackslash}p{#1}}
\setlist[itemize]{leftmargin=5mm,itemsep=1.2mm,topsep=2mm,parsep=0pt}

\begin{document}

%======================================================================
\panel{How the friction is added}

Garcia's force vector in $(X,\Theta,Y)$, Eq.~(6), one Coulomb term per physical actuator:

\begin{align*}
F_X&=(F_1+F_2)-\left[c_{\mathrm{c}1}\sign(\dot X_1)+c_{\mathrm{c}2}\sign(\dot X_2)\right]\\
F_\Theta&=\left\{(F_1-F_2)-\left[c_{\mathrm{c}1}\sign(\dot X_1)-c_{\mathrm{c}2}\sign(\dot X_2)\right]\right\}\tfrac{L_b}{2}\\
F_Y&=F_y-c_{\mathrm{c}y}\sign(\dot Y)
\end{align*}

\newpage
%======================================================================
\panel{Why the friction is kept sharp}

\begin{equation*}
|v_i|>0:\quad F_i=c_{\mathrm{c}i}\sign(v_i),
\qquad\qquad
v_i=0:\quad |F_i|\le c_{\mathrm{c}i}.
\end{equation*}
\vspace{-4mm}

At rest friction matches whatever is applied, so $c_{\mathrm{c}i}$ is the threshold to overcome
before the rail moves at all ($16.8$, $18.35$, $11.6$~N).

\begin{itemize}[itemsep=0.4mm,topsep=1mm]
\item \textbf{Keep the sharp flip.} A smoothed $\sign$ is just a steep viscous damper, and the
      baseline already has viscous damping: it would absorb the mismatch. The sharp flip is the
      part the baseline cannot represent.
\item \textbf{That forces a choice at $v=0$.} $\sign(0)$ has no value. Setting it to zero
      removes the threshold, so any force moves the rail. Instead we hold the rail and solve for
      its friction, capped at $\cc$.
\item \textbf{Still Garcia's law.} He writes $\sign$ and identifies at constant velocity, so he
      measures sliding and never treats $v=0$. The stick branch is the Filippov solution of
      his $\sign$ law, not a change to it.
\end{itemize}

\newpage
%======================================================================
\panel{How it enters our model}

Friction is subtracted from the applied force. The state-space representation is untouched:

\begin{align*}
\text{frictionless truth:}&\qquad \dot x=A(x)\,x+B(x)\,u\\
\text{with friction:}&\qquad \dot x=A(x)\,x+B(x)\bigl(u-PF\bigr)
\end{align*}

\begin{equation*}
v=P^{\T}\dot q=\col(\dot X_1,\dot X_2,\dot Y),
\qquad
P=\left[\begin{smallmatrix}1&1&0\\ L_b/2&-L_b/2&0\\ 0&0&1\end{smallmatrix}\right],
\qquad
\begin{aligned}
|v_i|&>v_\varepsilon: &&F_i=c_{\mathrm{c}i}\sign(v_i)\\
|v_i|&\le v_\varepsilon: &&F_i\ \text{solved, rail holds},\ |F_i|\le c_{\mathrm{c}i}
\end{aligned}
\end{equation*}
\vspace{-4mm}

$v_\varepsilon$: below it a rail counts as stopped.

\begin{itemize}[itemsep=0.4mm,topsep=1mm]
\item $M(Y,\delta_a)$, $C$, $K$ untouched, so $A$ and $B$ are the frictionless model's.
\item $v$ is a rail velocity, not a logical one. Same $P$ as the rest of the pipeline.
\item $F\in\R^3$: $\delta_a$ is an internal payload DOF, not a rail.
\item $\cc=0\Rightarrow F=0$, reproducing the frictionless function exactly.
\end{itemize}

\newpage
%======================================================================
\panel{Garcia's experiment, replicated}

His, Section~2.4, with $c_y,c_{\mathrm{g}i}$ the viscous coefficients:
\begin{align*}
\text{$Y$, with $X_1,X_2$ blocked:}&\qquad F_y=c_y\dot Y+c_{\mathrm{c}y}\sign(\dot Y)\\
\text{$X_1,X_2$ synchronised:}&\qquad F_i=c_{\mathrm{g}i}\dot X_i+c_{\mathrm{c}i}\sign(\dot X_i),
\quad i=1,2
\end{align*}
\vspace{-4mm}

Two runs: he reads each actuator force separately, so $c_{\mathrm{c}1}$ and $c_{\mathrm{c}2}$
come out directly.

Ours: the input is $u=\col(u_X,u_\Theta,u_Y)$, so the rails are mixed. Prescribe the velocity,
solve for the $u$ that holds it, fit $u$ on $v$ per sign branch.
\begin{align*}
\text{E1}\quad \dot Y=v:&\qquad u_Y=c_{\mathrm{c}y}\\
\text{E2}\quad \dot X_1=\dot X_2=v:&\qquad u_X=c_{\mathrm{c}1}+c_{\mathrm{c}2},
\qquad u_\Theta=\tfrac{L_b}{2}(c_{\mathrm{c}1}-c_{\mathrm{c}2})\\
\text{E3}\quad \dot X_1=v,\ \dot X_2=-v:&\qquad u_X=c_{\mathrm{c}1}-c_{\mathrm{c}2},
\qquad u_\Theta=\tfrac{L_b}{2}(c_{\mathrm{c}1}+c_{\mathrm{c}2})
\end{align*}
\vspace{-4mm}

\begin{itemize}[itemsep=0.4mm,topsep=1mm]
\item Three runs, because $u$ mixes the rails: $c_{\mathrm{c}1},c_{\mathrm{c}2}=$
      (E2 $\pm$ E3)$/2$ on $u_X$.
\item With $Y$ and $\delta_a$ held at zero, $M$ is constant, so $u$ is exactly affine in $v$
      per sign branch. Slope is viscous and inertial; the intercept is Coulomb.
\end{itemize}

\newpage
%======================================================================
\panel{What it returns}

\begin{tabular}{@{}lrr@{}}
\toprule
 & recovered [N] & Garcia [N] \\
\midrule
$c_{\mathrm{c}1}$ & $16.800000$ & $16.80$ \\
$c_{\mathrm{c}2}$ & $18.350000$ & $18.35$ \\
$c_{\mathrm{c}y}$ & $11.600000$ & $11.60$ \\
\bottomrule
\end{tabular}

\begin{itemize}
\item The fit is a perfect straight line, so the number we read off really is the friction.
\item Both rails came back with their own value, $16.8$ and $18.35$. Only possible if the
      friction sits on the rails.
\item A twist appeared that we never specified: geometry predicts
      $\tfrac{L_b}{2}(c_{\mathrm{c}1}-c_{\mathrm{c}2})=-0.5619$~Nm, the model returned
      $-0.561875$~Nm.
\end{itemize}

\newpage
%======================================================================
\panel{What had to be decided}

\begin{tabular}{@{}L{0.33\textwidth}L{0.55\textwidth}@{}}
\toprule
decision & why \\
\midrule
Friction in the truth only
& The augmentation needs something real to account for. Friction in the physics
block answers a different question. \\[2mm]
Garcia's $\cc$ as given, no sweep
& Fitting a physical constant to the outcome it explains is not a measurement. \\
\bottomrule
\end{tabular}

\end{document}
